\documentclass[output=paper,colorlinks,citecolor=brown]{langscibook} 

\bibliography{localbibliography}

%This is where you put the authors and their affiliations
\author{Michael Diercks\affiliation{Pomona College}\and Franny Brogan\affiliation{Pomona College}\lastand  Hazel Mitchley\affiliation{Rutgers University}}

%Insert your title here
\title{Field Methods Papers: Quick Reference Guide for LaTeX}  
\abstract{Here is where you write the abstract of the paper.
}

\input{FieldMethodsTemplatePackages.tex}
\input{FieldMethodsTemplateCommands.tex}

% \IfFileExists{../localcommands.tex}{%hack to check whether this is being compiled as part of a collection or standalone
%   \input{../localpackages}
%   \input{../localcommands} 
% \togglepaper[23]
% }{}

\begin{document}
\maketitle

%\tableofcontents 

\section{Introduction}

%Please note, this template is built on Language Science Press's LaTeX template from Overleaf (as of early 2020). It has been posted with their permission, but any errors should be attributed to Diercks and nobody else. If you intend to write for Language Science Press independently, please download their LaTeX templates directly from their website. http://langsci-press.org/

\textbf{Note}: This document is specifically designed to be read as a \href{https://www.overleaf.com/latex/templates/field-methods-papers-quick-reference-guide/cqbjxwnffvrm}{project on Overleaf}: double clicking on specific content on the right side of your screen will show you the source code to the left. So you can read the formatted document on the right of the screen, then look to the left to see the code that generated it.

\LaTeX{} is a typesetting program that takes a text document (with the file type .tex) and converts it to a typeset document that can be easily made into a PDF. Overleaf is a web-based app for writing in \LaTeX{}. In Overleaf, the .tex document is in the middle side of your screen; this `code' is compiled to generate typeset document is on the right, which can be downloaded as a PDF using the `Download PDF' button  at the top of the preview.

When writing a .tex document you directly annotate your text with commands that tell the program how to format the text. So if you want to write something in \textit{italics} you write \verb|\textit{italics}|, and if you want something \textbf{bolded} you write \verb|\textbf{bolded}|. \LaTeX{} then converts that `code' into what looks like italics and bold. In this Overleaf template, you can look at the .tex file itself to see what commands (`code') are used to create different kinds of formatting. 

%In the paragraph above, the command \verb|stuff| is from the package ``verbatim'' which prints whatever is written inside the |vertical lines| even if it is a command that usually does something in LaTeX. Just in case you were wondering what that command is.

Double click on this sentence in the preview and see where the cursor moved to in the .tex document to the left of this preview, to see what a ``comment'' is using the \% symbol. 

%As you can see, using the percentage symbol "%" creates a "comment," i.e. something that appears in the .tex document, but which does not compile, and therefore has no effect on the document being produced. It's a handy way to write yourself notes. This does mean that if you ever want to use a % symbol that is NOT for a comment, you have to add a backslash in front of it to tell LaTeX that you don't mean for it to be a comment, like we did in the paragraph above.

%If you want to comment out a bunch of stuff at once in Overleaf, highlight it and Press Cmd-/ (or, Ctrl-/ on a PC). This will but a % in front of every line.

\subsection{What is \LaTeX{} and why am I using it?}

These are very good questions: short story, our research team is publishing with \href{https://langsci-press.org}{Language Science Press} (LSP) for multiple reasons, chief of which is that it is Open Access, meaning that the books it publishes are available to download (for free) online. This is important because much academic work is hidden behind paywalls, and there are serious questions of inequity if our research on understudied/marginalized languages cannot be read by members of the communities that speak those languages. 

However, LSP requires submissions in \LaTeX{}, so we are writing in \LaTeX{}.

\subsection{What is this document?}

This is a Quick Reference Guide for writing papers in \LaTeX{} for Field Methods classes with Prof. Diercks at Pomona College, distributed as \href{https://www.overleaf.com/latex/templates/field-methods-papers-quick-reference-guide/cqbjxwnffvrm}{an Overleaf template}.

\begin{itemize}

\item This guide offers some basic introduction to \LaTeX{}, though it is not a full \LaTeX{} tutorial. For this, see \href{http://ptmartins.info/latex/}{Pedro Martin's tutorial}.

\item This guide provides examples of formatting you may need for a field methods paper, e.g. tables, numbered examples, etc. 

\item If you clicked on ``\href{https://www.overleaf.com/latex/templates/field-methods-papers-quick-reference-guide/cqbjxwnffvrm}{Open as a Template},'' a copy of this project is now in your Overleaf account (i.e. it belongs to you), and you can now edit this document to be an actual paper that you write. This means there is no setup required, you can just erase this text and start typing your own paper, or use this document to practice and see what changes your edits make.\footnote{This also means you will probably want to rename the document - you can do this by hovering your cursor over the title ``Field Methods Paper Quick Reference Guide' at the top of this page, and then clicking on the pencil icon that appears to the right of the title. You can also rename the FieldMethodsQuickReference.tex file in the leftmost panel by clicking on the down arrow.} 

\item Designing a \LaTeX{} project from scratch can be difficult; the point of this template is that this is already done for you, you can just start writing. The \href{https://www.overleaf.com/latex/templates/field-methods-paper-template/pqrxmrsptqcv}{Field Methods Paper Template} offers the same advantages without this (long) explanatory document.

\item Don't worry about breaking anything - you can always open the template link again from Overleaf's website to see this initial version again. So edit away!

%Mess with the text in the .tex document right here and see what happens when you Recompile (the green button at the top left of the PDF viewer).

\end{itemize}

\section{Some basic things you will want to know}

This section \textit{very} quickly outlines things you may want to know how to do to write field methods papers. 

\subsection{Useful commands for text formatting, some that we built for you}
 \label{sec:diercks:UsefulShortcuts}

The chart below is structured with the most general sorts of text formatting first, with linguistics-specific symbols/formatting at the bottom.

\begin{table}
\caption{Selected commands for common ACAL formatting needs}
\label{tab:diercks:CommandsTable}

\begin{tabular}[t]{l l l}
\lsptoprule		
\textbf{Symbol/Annotation} & \textbf{Example} & \textbf{Code} \\ \midrule
		Ellipsis & \dots & \verb|\dots| \\
		Subscript & NP\subs{i} & \verb|NP\subs{i}| \\
		Superscript & NP\supers{i} & \verb|NP\supers{i}|\\ 
		Bold & \textbf{bold} & \verb|\textbf{bold}| \\
		Italic & \textit{italic} & \verb|\textit{italic}| \\
		Small Caps & \textsc{small caps} & \verb|\textsc{small caps}| \\
		Strikeout & \sout{strikeout} & \verb|\sout{strikeout}| \\
		Underline & \uline{underline} & \verb|\uline{underline}| \\
		circle something in text & \circled{something} & \verb|\circled{something}| \\
		Highlight something & \hl{something} & \verb|\hl{something}| \\
		Null & \nothing  & \verb|\nothing|\\
	\lspbottomrule
	\end{tabular}
\end{table}

\begin{itemize}

\item If you want to make a bulleted list, look at how this list is formatted in the .tex document using the ``itemize'' environment. Notice \verb|\begin{itemize}| which starts the bulleted list, \verb|\item| which introduces a bullet, and \verb|\end{itemize}|, which ends the bulleted list.

\item As you've already seen if you are paying attention to the .tex document on the left of your screen, sections, subsections, and sub-subsection are formatted with the commands \verb|\section{}|, \verb|\subsection{}|, and \verb|\subsubsection{}|, respectively.

\item Look at the .tex document to see how we bolded \textbf{this text} (and Overleaf has a shortcut to make it easy, Cmd-B on Macs, Ctrl-B on PCs). \textit{Similarly for italics}, Overleaf provides a shortcut (Cmd-I on Macs, Ctrl-I on PCs).

\item Write footnotes like this.\footnote{Hey, look at me, I'm a footnote.}

\end{itemize}

\subsection{How do I do X?}

You likely will wonder how to do something that seems mysterious in \LaTeX{} but that is so simple in other applications. Google is your friend. Type ``bulleted list in latex" into Google and you will quickly find instructions for a bulleted list, and you can use this method to find out how to do almost anything you need to do. If you still cannot figure out how to solve your \LaTeX{} puzzle, talk to your classmates, your TA, or Professor Diercks.  

\subsection{Heads up! Some potential pitfalls}

For all of \LaTeX{}'s conveniences, there are some annoyances that are puzzling until you figure them out. 

\begin{itemize}

    \item Quotation marks are easy to get wrong. You want them to display ``like this'' and not "like this," so use the grave accent symbol above your tab key to type front quotes - use one for a single quote `like this' and two for quotation marks ``like this.'' For reference, the grave accent is the one that appears on the mid vowel here: \`o.

    \item Because a character like the percentage symbol \% actually means something in \LaTeX{} (it creates a comment), you can't just use a percentage symbol on its own to mean a percentage symbol. Instead, you have to put a backslash in front of it, like we did in in the previous sentence (check out the .tex document, and also \href{http://ptmartins.info/latex/#special-characters}{this note} about other similar characters).
    
    \item It is very easy to forget to put the end bracket on a command. So you might write this \verb|\hl{highlight this| instead of \hl{highlight this} (see the .tex) and your document will show an error. Thankfully Overleaf has warnings that pop up as you type if you have left a bracket open.
    
    \item Because it's possible to write code that won't compile, it is recommended to compile your document frequently while writing. This way, when you make a mistake you won't have pages and pages to sort through to find where the mistake is. The keystroke Cmd-S (on Mac) and Ctrl-S (on PC) will instruct Overleaf to Recompile, we suggest making it a habit to do this frequently while you write. %This makes things a lot easier, bc trying to find the broken piece of code somewhere inside your whole paper can be a real pain.

\end{itemize}

\section{Numbered examples} \label{sec:diercks:NumberedExs}

It is easy to create consecutively numbered examples: 

%\ea Watch the numbers change when you uncomment this and recompile!  \z 

\ea 
This example is auto-numbered - if you uncomment the example above that is in green in the .tex document (by deleting the \% before the example) and recompile, the numbering will change.
\z 

\noindent The \texttt{gb4e} package that Language Science Press uses for numbered examples allows for inter-linear glossing and translations: 

\ea 
%\langinfo{Language Name here}{}{source here} %this line is optional
\gll this is a language example \\
this is the morphological gloss \\
\glt `This is the translation.' 
\z 

\noindent  \texttt{gb4e} automatically aligns glosses with the language example (each white space means a different word). This is useful for non-English examples like the Lubukusu sentence in \REF{ex:diercks:BukusuNegation}: 

\ea \label{ex:diercks:BukusuNegation}
%\langinfo{Lubukusu}{}{Justine Sikuku, pc}
\gll Wekesa se-a-la-ba a-kula ka-ma-indi ta. \\
Wekesa \Neg-1\Sm-\Fut-be 1\Sm-buy 6-6-maize \Neg{} \\ 
\glt `Wekesa will not be buying maize.'
\z

\noindent Some things to notice in the .tex document for the examples above - an example is started with the \verb|\ea| and ended with \verb|\z| - omitting either of these will generate an error. The language and gloss lines are started with the command \verb|\gll|, with line breaks (\verb|\\|) ending both the language line and the gloss line. The translation line is introduced with the command \verb|\glt|, and notice that there is no line break at the end of the translation line. If you want to label an example as coming from a specific language, Language Science Press uses the \verb|\langinfo{}{}{}| command that is commented out in example \REF{ex:diercks:BukusuNegation}.

If you want multiple data examples to appear in one numbered example, you use the xlist environment to create a sublist, which you can see in the .tex document by double clicking on the example below in the preview: 

\ea \label{ex:diercks:XlistExample}
\begin{xlist}

\ex This is the first example. 

\ex \label{ex:diercks:SwahiliXlist}
\gll hii ni m-fano w-a pili, kwa Ki-swahili \\
this is 3-example 3-\Assoc{} second, of 7-Swahili \\
\glt `This is a second example, in Swahili.' 

\end{xlist}

\z 

\noindent Again, some details to notice in \REF{ex:diercks:XlistExample}: first, notice the commands \verb|\begin{xlist}| and \verb|\end{xlist}| that begin and end the sublist inside the example. It is important that the sublist be generated \textit{inside} the numbered example, so \verb|\begin{xlist}| comes after \verb|\ea| and \verb|\end{xlist}| comes before \verb|\z|. Each example inside the sublist is introduced with the command \verb|\ex|, and each of these examples uses normal gb4e formatting. So you can see that the code to generate \REF{ex:diercks:SwahiliXlist} is the same as was used to generate \REF{ex:diercks:BukusuNegation}.

In many linguistics publications, \textit{all} diagrams and data are presented as numbered examples (e.g. trees, tableaux, charts, sentences with interlinear glosses, data sets of words). However, in Language Science Press tables are created as Tables instead of numbered examples (see \sectref{sec:diercks:BasicTables} and \sectref{sec:diercks:AdvTabs}).\footnote{Syntactic trees are formatted as figures with LSP, but we won't need trees for field methods papers.} 

\section{Using cross-references} \label{sec:diercks:CrossReferences}

It is useful to be able to refer your reader to different sections/examples in your paper: for example, we may want to tell you that examples \REF{ex:diercks:BukusuNegation} and \REF{ex:diercks:XlistExample} are built using \texttt{gb4e}, or that if you want to learn about fonts you should read \sectref{sec:diercks:fonts}. There are two components of a cross-reference in text: 

\begin{enumerate}    

\item You must label a particular section or example that you intend to refer to using the command \verb|\label{LabelName}|, where ``LabelName'' has some predetermined portions as well as a portion where you use some text description to describe the thing you are labeling:\footnote{These commands are not standard for \LaTeX{}, they are specifically part of LSP's template.} 

\begin{itemize}

\item Section label: \verb|\label{sec:author:yourlabel}|

\item Table label: \verb|\label{tab:author:yourlabel}|

\item Figure label: \verb|\label{fig:author:yourlabel}|

\item Example label: \verb|\label{ex:author:yourlabel}|

\end{itemize}
    
    
\item You must then use one of these commands to create the cross-reference. 

\begin{itemize}

\item Section reference (will read `\S \#’ in output):\\ \verb|\sectref{sec:author:yourlabel}|

\item Table reference (will read `Table \#’ in output):\\ \verb|\tabref{tab:author:yourlabel}|

\item Figure reference (will read `Figure \#’ in output): \\ \verb|\figref{fig:author:yourlabel}|

\item Example reference with parens — note caps (will read `(\#)’ in output): \\ \verb|\REF{ex:author:yourlabel}| 

\item Example reference without parens — note caps (will read `\#’ in output): \\ \verb|\ref{ex:author:yourlabel}|

\end{itemize}

\item Some examples:\footnote{Look at the .tex on the left to see how to write the commands, and click the hyperlinks in the cross-references to take you to the portion of the document where the label is created.}

\begin{itemize}

\item the Swahili example above can be referenced as \REF{ex:diercks:SwahiliXlist}. If I wanted to refer to it without parentheses, I can use this command: \ref{ex:diercks:SwahiliXlist}.

\item The aforementioned examples occur in \sectref{sec:diercks:NumberedExs}, which describes how to use numbered examples.

\item  In \sectref{sec:diercks:BasicTables},  \tabref{tab:diercks:SampleTableCenter} shows how to center-align text in cells of a table.

\item These are some selected examples - throughout the entire document you will see these kinds of cross-references being utilized, so you can look at numerous examples if you so wish.

\end{itemize}

\end{enumerate}

Notably, you can only use any of the reference commands (e.g. the command for numbered examples, \verb|\REF{LabelName}|) with labels that already exist somewhere in your paper - attempting to build a cross-reference to a label that does not exist will not stop your document from compiling, but it will result in a Yellow Warning from Overleaf next to the Recompile button, and it will result in question marks where the cross-reference ought to be, like this: \REF{ToNothing}. %If you change the reference in the preceding sentence to \REF{ex:diercks:SwahiliXlist} instead of ``ToNothing,'' you should see the Yellow Warnings next to the Recompile button reduce in number by two, and the cross-reference will successfully compile. 

These cross-references are auto-generated every time you compile your document, so as your section/example numbers change as you write your paper, the cross-references update themselves. 

\section{Using IPA symbols} \label{sec:diercks:fonts}

\ea 
ðɪs ɪz haʊ tu raɪt ɪn ʔɑɪpijeɪ. (This is how to write in IPA.)
\z 

\noindent This template is built on Language Science Press's template, which which allows you to enter IPA symbols directly into your .tex document and it appropriately typesets it when creating a PDF. This doesn't tell you how to type the IPA symbols in the first place, though. This requires a non-\LaTeX{} solution.

\begin{itemize}

    \item \textbf{Recommended}: \href{https://scripts.sil.org/cms/scripts/page.php?site_id=nrsi&id=uniipakeyboard}{SIL IPA keyboard} is downloadable on Mac and PC and gives you keystrokes for inserting IPA characters.\footnote{FYI there's no way to type a `\textrhookschwa' via this keyboard (as far as we know), so utilize one of the methods below.}
    
    \item \href{https://www.blugs.com/IPA/}{IPA Palette}
    
    \item \href{https://ipa.typeit.org}{Online IPA keyboard} where you can type symbols and then copy/paste them into your document
    
    \item The web app \href{http://detexify.kirelabs.org/classify.html}{detexify} is a handy tool that allows you to draw the symbol you want and it shows you what packages you need for that symbol, and what commands create it in \LaTeX{}.
    
    \item The web app \href{https://shapecatcher.com}{shapecatcher} does something similar; you draw the symbol you want and it identifies possible matches and their unicode identifiers.
    
    

\end{itemize}

%These are TIPA instructions for those who want them. But I (MJKD) am thinking we can leave this out altogether, bc I think it won't be necessary?? 

% \noindent If you \textbf{do} want a \LaTeX{}-based solution for creating IPA symbols (i.e. if you don't want to use the keyboards above), the \texttt{TIPA} package is the most commonly-used system for processing IPA symbols in \LaTeX{}. (\href{https://www.tug.org/TUGboat/tb17-2/tb51rei.pdf}{This document} is a good resource for the macro names (i.e. commands to draw a symbol) and shortcut characters you'll need if you want to use \texttt{TIPA}.) Crucially, \texttt{TIPA} phonetic symbols can be inputted in the following two ways (look at the .tex document to see how these are coded):
% \begin{enumerate}
%     \item Input macro names in the normal text environment, e.g., \dh\textsci s \textsci z ha\textupsilon\ tu ra\textsci t \textsci n \textglotstop\textscripta pije\textsci.
%     \item Input macro names or \textit{shortcut characters} within the following groups or environment:
%     \begin{itemize}
%         \item \verb|\textipa{...}|, e.g., 
%             \textipa{DIs Iz haU tu raIt In PApijeI}.
%         \item \verb|{\tipaencoding ...}|, e.g., 
%             {\tipaencoding DIs Iz haU tu raIt In PApijeI}.
%         \item \verb|\begin{IPA} ... \end{IPA}|, e.g.,
%             \begin{IPA}
%             DIs Iz haU tu raIt In PApijeI.
%             \end{IPA}
%     \end{itemize}
% \end{enumerate}

% \noindent In our estimation, \texttt{TIPA} is functional if you have to write the occasional IPA character. But the learning curve for \texttt{TIPA} is not much different than the learning curve for the IPA keyboard (and it is harder to read the code in a document coded with \texttt{TIPA}, as you can see in the examples in the list above), so if you are frequently using IPA symbols, you will likely find it easier to learn the IPA keyboard linked to above.


\section{Citing References}

This section is not a proper introduction to all of how citations and bibliographies work in \LaTeX{}. Rather, we attempt to give you the basics you need to know in order to generate citations and a bibliography for a field methods paper. 

\subsection{How bibliographies work}

You will notice on the left side of the Overleaf display there is a list of files. The file titled ``localbibliography.bib'' is the bibliography file, and contains references that can be cited in your document. Click on it to see what a \LaTeX{} reference list looks like, and then click on ``FieldMethodsQuickReference.tex'' to come back here. 

As you write, you can type commands into your text which will do two things: first, it will generate citations in the text to the paper you cite; second, when your document compiles it will add that reference to your bibliography at the end of your paper. So if you want to make reference to \citet{BresnanMchombo1987ChichewaSmOm}, enter the citation command (see the .tex document), which references the file bibliographic entry in ``localbibliography.bib'' and generates the appropriate citation here in this paragraph, and adds the \citeauthor{BresnanMchombo1987ChichewaSmOm} paper to the bibliography at the end of this document.

The subsections that follow explain in more detail how to do this in your own paper. 

\subsection{Building a bibliography in a .bib file}

There are many ways to build your reference list: we outline two major ones here.\footnote{If you consistently write in \LaTeX{} you eventually would build a large .bib file that you use on every project; our goal in this explanation is mainly to help novices learn enough to write their current paper.}

First, if you have already created bibliographic references, you can use \citet{Nordhoff2018} (which is located \href{http://glottotopia.org/doc2tex/doc2bib}{here}) to convert those references to a format appropriate for \LaTeX{}. You can then copy and paste those references into ``localbibliography.bib'' which makes them available to cite in your paper.

Second, you can download a reference manager, which produces reference lists that look like what you see in ``localbibliography.bib'' but which gives you a user-friendly interface. If you are using a Mac, \href{https://bibdesk.sourceforge.io}{BibDesk} is an excellent application that will manage your citations - a brief introductory video is \href{https://www.youtube.com/watch?v=sOxf1bKAxqQ}{available on Youtube}. If you are using a PC, we recommend using \href{http://www.jabref.org}{Jabref}, which has the same functionality. Those reference managers will allow you to build a .bib file, which can be copied/pasted into ``localbibliography.bib.'' 

\subsection{What is a cite key?}

If you use the \href{http://glottotopia.org/doc2tex/doc2bib}{Bib converter}, you will notice that the first line of a bibliographic entry is auto-generated as ``AuthorDate'' which is ``Bloomfield1925'' in the example below (which is the first example in the \href{http://glottotopia.org/doc2tex/doc2bib}{Bib converter} sample). 

\indent @article\{\circled{Bloomfield1925}, \\
\indent \indent author = {Bloomfield, Leonard}, \\
\indent\indent journal = {Language}, \\
\indent\indent number = {4},\\
\indent\indent pages = {130--156},\\
\indent\indent title = {On the sound-system of central {Algonquian}},\\
\indent\indent volume = {1},\\
\indent\indent year = {1925}\\
\indent\indent\} \\

\noindent This circled text above is the ``cite key,'' which is an identifier that you use to cite that reference (this can in fact be any string of text that you choose).

\subsection{Citing sources}
%This section was authored by Hazel Mitchley, originally as part of an early draft of the ACAL Proceedings LaTeX template.

The two most important citation commands are \verb|\citet{}| and \verb|\citep{}|. The cite key of the reference goes between the curly brackets.
 
 \begin{itemize}
 \item For `Author (Date)' format, use \verb|\citet{}|. \\
  e.g.: writing \verb|\citet{Lahiri2000}| produces \citet{Lahiri2000}
 \item For a parenthetical citation, use \verb|\citep{}|.\\
 e.g.: writing \verb|\citep{Lahiri2000}| produces \citep{Lahiri2000}
 \end{itemize}

\noindent Both of these commands can be modified to add extra information.
 
 \begin{itemize}
 \item To add a page number after the year, put the page number in square brackets \textit{before} your curly brackets.\\
 e.g.: \verb|\citep[87]{Doke1923}| produces \citep[87]{Doke1923} 
 \item To add information before the citation, put this in square brackets before the brackets with the page number.\\
  e.g.: \verb|\citep[Zulu,][87]{Doke1923}| produces \citep[Zulu,][87]{Doke1923} 
  \item If you want information before the citation, but without a page number, just leave the page number blank.\\ 
  e.g. \verb|\citep[Zulu,][]{Doke1923}| produces \citep[Zulu,][]{Doke1923}
  \item Multiple cite keys separated by commas will produce multiple citations.\\
  e.g. \verb|\citep{Doke1923,Lahiri2000}| produces \citep{Doke1923,Lahiri2000}
 \end{itemize}
 

\noindent You can check out this \href{http://merkel.texture.rocks/Latex/natbib.php}{helpful reference sheet} that gives many details about citation commands for any additional variants that you want. %LSP uses biblatex, not natbib, but their template (and this one) are designed so that the natbib commands work. 

If you need help getting your references set up, please contact Prof. Diercks, or the course TA.

\section{Tables}

\subsection{Basic tables} \label{sec:diercks:BasicTables}

Tables are very useful for formatting data and examples, but counter-intuitive to produce in \LaTeX{} as compared to in familiar word processors. This section lays out the basics; see \sectref{sec:diercks:AdvTabs} for additional formatting tricks. You can also check out \href{https://en.wikibooks.org/wiki/LaTeX/Tables}{this useful wiki} on how to construct \LaTeX{} tables if you want some addition information. 

The \texttt{tabular} environment is used to typeset tables. By default, \LaTeX{} tables are drawn without any vertical or horizontal lines and column width is predetermined; this means that any settings beyond these defaults must be defined by you. Per LSP's guidelines, tables should have vertical borders in field methods papers. Let's take a look at some simple examples below. Please note, tables are moved around by \LaTeX{} in the typeset document on the right of your screen to where they typeset most naturally, so you may have to look to the next page for a table. Also, please pay attention to the commented-out text that appears in green in the .tex document around each table's code, because we have explained what each part of the code for the table does.   %If for some reason you view vertical borders as a necessity, please get in contact with your volume's editors rather than submitting a paper that includes them.


Let's start with a basic 3x3 table with single horizontal borders. Take a look at the .tex doc to see how borders and column alignment are specified (i.e. whether your text is aligned on the left, right, or in the center of the column) for \tabref{tab:diercks:FirstSampleTable} below. 


\begin{table}

\caption{A descriptive caption here}
\label{tab:diercks:FirstSampleTable}

\begin{tabular}{lll} %the three `l's' mean there are three columns, and that they are left-aligned (hence `l'). 
\lsptoprule %This (\lsptoprule) is a required piece of formatting for tables in LSP volumes (and, therefore, in Field Methods papers) 
cell 1 & cell 2 & cell 3 \\ 
\hline %comment this out to remove horizontal border
cell 4 & cell 5 & cell 6 \\
\hline %comment this out to remove horizontal border
cell 7 & cell 8 & cell 9 \\
\lspbottomrule %This (\lspbottomrule) is a required piece of formatting for tables in LSP volumes (and, therefore, in Field Methods papers) 
\end{tabular}

\end{table}

\noindent If you want to remove the horizontal borders, try commenting out the commands \verb|\hline| in the table above; this will remove the horizontal borders.

%Note that the contents of cells in a given row are separated by '&'. You should have the same number of '&' signs in every row, even if the contents of a given cell are blank (i.e. a blank cell will look like this: & &). Also note that new rows are created using '\\'. This is a more general LaTeX command you can use for manually starting a new line.

\tabref{tab:diercks:SampleTableHeaders} is a table that's identical to the one above, but with column headers separated by a double horizontal line. Take a look at the code in the .tex document to see how it's different from the code that generated the table above. %\verb|You can generate as many horizontal lines as you'd like by duplicating the \hline command.| 

\begin{table}

  \caption{This table separates headers with a double line}
  \label{tab:diercks:SampleTableHeaders}
  
\begin{tabular}{lll}
\lsptoprule %This (\lsptoprule) is a required piece of formatting for tables in LSP volumes (and, therefore, in Field Methods papers) 
Header 1 & Header 2 & Header 3 \\
\hline\hline
cell 1 & cell 2 & cell 3 \\
\hline
cell 4 & cell 5 & cell 6 \\
\hline
cell 7 & cell 8 & cell 9 \\
\lspbottomrule %This (\lspbottomrule) is a required piece of formatting for tables in LSP volumes (and, therefore, in Field Methods papers) 


\end{tabular}
\end{table}

\noindent Finally, \tabref{tab:diercks:SampleTableCenter} is a table identical to the one above, but with text center-aligned.

\begin{table}

  \caption{This table has center-aligned text}
  \label{tab:diercks:SampleTableCenter}

\begin{tabular}{ccc}
\lsptoprule %This (\lsptoprule) is a required piece of formatting for tables in LSP volumes (and, therefore, in Field Methods papers) 
Header 1 & Header 2 & Header 3 \\
\hline\hline
cell 1 & cell 2 & cell 3 \\
\hline
cell 4 & cell 5 & cell 6 \\
\hline
cell 7 & cell 8 & cell 9 \\
\lspbottomrule %This (\lspbottomrule) is a required piece of formatting for tables in LSP volumes (and, therefore, in Field Methods papers) 

\end{tabular}
\end{table}

\subsection{Advanced tables}\label{sec:diercks:AdvTabs}
In \sectref{sec:diercks:BasicTables}, we showed you how to generate basic tables in \LaTeX{}. This section provides guidelines for a few (mostly stylistic) modifications that may be useful for you. For a comprehensive ``how-to" guide, see Overleaf's fantastic \href{https://www.overleaf.com/learn/latex/Tables}{tables tutorial}, or \href{https://en.wikibooks.org/wiki/LaTeX/Tables}{this wiki}. 

\subsubsection{Controlling basic table formatting}

%For those of us who just can't let something else be in control

\begin{itemize}
   
\item \textbf{Manually define column width}. Thus far, our column width has been determined by the widest cell. By using \verb|p{}| when specifying your parameter, you can manually define the width of each individual column, inducing line breaks and/or justified columns.

\begin{table} 
\caption{\large{Now my columns are wider}}
\centering
\begin{tabular}{p{1in}p{1in}p{1in}} %this specifies 1-inch columns
\hline
\textbf{Header 1} & \textbf{Header 2} & \textbf{Header 3} \\
\hline\hline
cell 1 & cell 2 & cell 3 \\
\hline
cell 4 & cell 5 & cell 6 \\
\hline
cell 7 & cell 8 & cell 9 \\
\hline
\end{tabular}
\end{table}

\begin{table}[h!] 
\caption{\large{Now my columns are narrower}}
\centering
\begin{tabular}{p{1cm}p{1cm}p{1cm}} %this specifies 1cm columns
\hline
\textbf{Header 1} & \textbf{Header 2} & \textbf{Header 3} \\
\hline\hline
cell 1 & cell 2 & cell 3 \\
\hline
cell 4 & cell 5 & cell 6 \\
\hline
cell 7 & cell 8 & cell 9 \\
\hline
\end{tabular}
\end{table}

%This is commented out, since LSP does not allow vertical borders in tables

% \item \textbf{Control vertical borders}. Simply add/remove `|` symbols from your parameter where you don't want vertical borders:

% \begin{table}[h!] 
% \caption*{\large{I'm a table with (vertical) borders}} 
% \centering
% \begin{tabular}{c|c|c} %Notice, the vertical lines are here, which tells \LaTeX{} to insert vertical borders 
% \hline
% \textbf{Header 1} & \textbf{Header 2} & \textbf{Header 3} \\
% \hline\hline
% cell 1 & cell 2 & cell 3 \\
% \hline
% cell 4 & cell 5 & cell 6 \\
% \hline
% cell 7 & cell 8 & cell 9 \\
% \hline
% \end{tabular}
% \end{table}

\item \textbf{Control horizontal borders}. Simply remove any \verb|\hline| commands where you don't want horizontal borders:

\begin{table}[h!] 
\caption{\large{I'm a table without any borders at all}} 
\centering
\begin{tabular}{c c c}
\textbf{Header 1} & \textbf{Header 2} & \textbf{Header 3} \\
cell 1 & cell 2 & cell 3 \\
cell 4 & cell 5 & cell 6 \\
cell 7 & cell 8 & cell 9 \\%No \hline commands between lines anymore
\end{tabular}
\end{table}

\end{itemize}

\subsubsection{Combining columns and rows in tables}
Columns and rows can be combined in a bigger cell using the \verb|\multicolumn{}{}{}| and \verb|\multirow{}{}{}|commands.

\noindent The \verb|\multicolumn{}{}{}| command takes three arguments:
\begin{enumerate}
    \item The number of columns to be combined 
    \item Delimiters and alignment of the resulting cell - i.e. do you want borders, and whether it's left-, center-, or right-aligned
    \item Text to be displayed inside the cell 
\end{enumerate}

In \tabref{tab:diercks:multicolumn} we see that the columns making up the top row have been combined into one using the \verb|\multicolumn{}{}{}| command.

\begin{table} 
\caption{Using multicolumn{}{}{}}
\label{tab:diercks:multicolumn}

\begin{tabular}{p{3cm}p{3cm}p{3cm}  }
\lsptoprule %This (\lsptoprule) is a required piece of formatting for tables in LSP volumes (and, therefore, in Field Methods papers) 

 \multicolumn{3}{c}{Pomona/Pitzer Linguistics Professors} \\
 \hline\hline
 Name & Initials & Name in IPA \\
 \hline
 Carmen & CF & \textipa{"kA\*r.m@n}\\
 Mary & MP & \textipa{"mE.\*ri}\\
 Mike & MD & \textipa{"maIk} \\
 Nicole & NH  & \textipa{n@."kol}\\
 Robin & RM & \textipa{"\*rA.b@n}\\
 Franny & FB & \textipa{"f\*r\ae.ni}\\
 \lspbottomrule %This (\lspbottomrule) is a required piece of formatting for tables in LSP volumes (and, therefore, in Field Methods papers) 
\end{tabular}
\end{table}

\noindent The \verb|\multirow{}{}{}| command also takes three arguments, but note that the second one is different than for \verb|\multicolumn{}{}{}|:
\begin{enumerate}
    \item The number of rows to be combined
    \item The width of the column 
    \item Text to be displayed inside the cell
\end{enumerate}

In the \tabref{tab:diercks:multirow}, we see that the rows in the first column have been combined:

\begin{table} 

\caption{Using multirow{}{}{}}
\label{tab:diercks:multirow}

\begin{tabular}{ cc }
\lsptoprule %This (\lsptoprule) is a required piece of formatting for tables in LSP volumes (and, therefore, in Field Methods papers) 

 Subfield & Professor(s) \\
 \hline\hline
 \multirow{4}{1in}{Sociolinguistics} & Carmen \\ 
 & Nicole\\
 & Franny\\
 & Robin \\
 \hline
 \multirow{2}{1in}{Phonology} & Mary \\
 & Franny\\
 \hline
 \multirow{2}{1in}{Phonetics} & Nicole \\
 & Franny\\
 \hline
 \multirow{2}{1in}{Syntax} & Mike \\ 
 & Robin \\
 \hline
 \multirow{2}{1in}{Fieldwork} & Mary \\ 
 & Mike \\
 \hline
 Psycholinguistics & Robin\\
 \lspbottomrule %This (\lspbottomrule) is a required piece of formatting for tables in LSP volumes (and, therefore, in Field Methods papers) 
\end{tabular}
\end{table} 

Notice in the code above that for each multi-row cell, it must be paired with the equivalent number of rows in any other columns. %So, the `Sociolinguistics' cell is 3 rows tall, so there must be three rows paired with it (here, all have text: Carmen, Nicole, Franny). 

\section{Conclusion}

This guide highlights the main kinds of formatting you will need for Field Methods papers, but you may well run into formatting questions not discussed here. There are a host of relevant resources online, and if you hit a confusing puzzle you can almost always google ``my problem latex'' and find relevant articles on Stackexchange and other similar sites. This guide is not meant to be comprehensive at all, but rather is meant to introduce novices to the central issues. If you have questions, please contact Professor Diercks or your TA.

\section*{Abbreviations}

To be added later.

%Abbreviations should be included in this section. It is a required part of volumes with Language Science Press that an abbreviations list be included, so please do not remove this section, but for Field Methods you don't have to write anything here.

\section*{Acknowledgements}

Acknowledgements should be listed here.

\noindent This \href{https://www.overleaf.com/latex/templates/field-methods-papers-quick-reference-guide/cqbjxwnffvrm}{Quick Reference Guide} and the \href{https://www.overleaf.com/latex/templates/field-methods-paper-template/pqrxmrsptqcv}{corresponding Paper Template} borrow directly (with permission) from corresponding documents written by Michael Diercks and Franny Brogan for linguistics students at Pomona College, and from the in-development \LaTeX{} template for proceedings papers for ACAL (Annual Conference on African Linguistics). These templates are built on Language Science Press's Overleaf template (for eventual publishing with LSP), and are published on Overleaf with LSP permission. Our thanks to Sebastian Nordhoff for his feedback on earlier versions of this template. Diercks' original \LaTeX{} teaching materials were created with assistance from Claire Halpert, Nico Baier, Jason Zentz, Maddy Bossi, and Mica Clausen. Our gratitude also to the ACAL \LaTeX{} committee for their input to the ACAL template, which also found its way here (especially Peter Jenks, Ken Steimel, and Matthew Faytak). 

\printbibliography[heading=subbibliography,notkeyword=this]

\end{document}
